\documentclass[10pt,letterpaper]{article}
\usepackage[top=0.85in,left=2.75in,footskip=0.75in]{geometry}
\usepackage{amsmath,amssymb}
\usepackage{changepage}
\usepackage[utf8x]{inputenc}
\usepackage{textcomp,marvosym}
\usepackage{cite}
\usepackage{nameref,hyperref}
\usepackage[right]{lineno}
\usepackage{microtype}
\DisableLigatures[f]{encoding = *, family = * }
\usepackage[table]{xcolor}
\usepackage{array}
\usepackage{listings}
\usepackage{graphicx}
\usepackage{float}
\usepackage{booktabs}
\usepackage{multirow}

\linenumbers

\title{MedImages.jl: A High-Performance, Differentiable Julia Framework for Medical Image Processing and Physics-Integrated Dosimetry}

\author{Jakub Mitura$^{1*}$, Divyansh Goyal$^{1}$, Joanna Wybranska$^{1}$, Michael C. Kreissl$^{1}$}

\date{}

\begin{document}

\maketitle

\noindent
$^{1}$ Division of Nuclear Medicine, Department of Radiology \& Nuclear Medicine, Faculty of Medicine, Otto von Guericke University Magdeburg, 39120 Magdeburg, Germany \\
$^{*}$ Corresponding author: jakub.mitura@med.ovgu.de

\section*{Abstract}
\textbf{Background and Objectives:} Quantitative medical imaging requires software architectures that ensure high-throughput computational performance and strict mathematical coupling between voxel data and spatial metadata. We introduce MedImages.jl, a native Julia ecosystem for differentiable 3D/4D medical image analysis that addresses the "two-language problem" by enforcing zero-cost runtime metadata abstractions via Julia's type system. 

\textbf{Methods:} MedImages.jl binds voxel data to spatial metadata (origin, spacing, direction) in a strictly typed struct. We validated the framework through four critical axes: (1) biobank-scale throughput using HDF5, (2) GPU-accelerated spatial kernels, (3) end-to-end differentiability via SciML, and (4) clinical metadata consistency. The primary validation is a novel Universal Differential Equation (UDE) dosimetry pipeline for $^{177}$Lu-PSMA.

\textbf{Results:} MedImages.jl achieved a 7.2$\times$ speedup over Python-based caching in a 100-subject benchmark. The UDE model achieved a Pearson correlation of \textbf{0.957} against Monte Carlo simulations ($N=48$ patches from 12 patients, $p < 0.001$), outperforming established clinical analytical standards and deep-learning baselines. Computational benchmarks prove that the Julia implementation is \textbf{10.5--13.6$\times$ faster} than established Python SciML frameworks (PyTorch/torchdiffeq and JAX/Diffrax).

\textbf{Conclusions:} By unifying spatial operations and metadata management within a single differentiable language, MedImages.jl facilitates high-fidelity clinical research without the ergonomic and performance overhead of hybrid frameworks.

\textbf{Keywords:} Medical Imaging, Julia, Scientific Machine Learning, Nuclear Medicine, Differentiable Programming, Spatial Metadata

\section{Introduction}
Medical imaging suffers from the "two-language problem": researchers prototype in Python but rely on opaque C++ kernels (SimpleITK, ITK) for performance. This divide introduces significant limitations in metadata preservation, differentiability, and computational efficiency, particularly in the demanding domain of nuclear medicine and radioligand therapy ($^{177}$Lu-PSMA).

The development of MedImages.jl is driven by the need to address four converging challenges in contemporary medical imaging pipelines.

\paragraph{Challenge 1: The Volume of Modern Biobank Datasets.}
A single whole-body PET/CT examination involves high-resolution 3D volumes that must be co-registered and resampled. At biobank scale, preprocessing thousands of subjects becomes a dominant bottleneck. Traditional hybrid environments suffer from significant I/O overhead and metadata synchronization delays when preparing multimodal datasets.

\paragraph{Challenge 2: Execution Speed and the "Two-Language" Barrier.}
Researchers often prototype in high-level languages like Python but must rely on opaque C++ backends for execution. Julia addresses this by compiling directly to efficient machine code via LLVM JIT, allowing spatial transformations to be executed rapidly across both CPU and GPU backends using native implementations.

\paragraph{Challenge 3: Differentiability and the Fragmented SciML Landscape.}
Emerging Scientific Machine Learning (SciML) paradigms require pipelines where every step is differentiable. While Python offers powerful differentiable libraries (PyTorch, JAX), they often function as "walled gardens." Julia allows native code to seamlessly interoperate through multiple dispatch, enabling end-to-end differentiability across complex differential equation solvers and GPU kernels.

\paragraph{Challenge 4: Metadata Management and Physical Interpretation.}
Unlike standard computer vision, medical images possess a strict "physical interpretation." A medical image is a grid of points occupying a physical region defined by origin, spacing, and direction. The transition to raw numerical arrays for deep learning often strips this spatial context, risking "metadata drift"—a dangerous misalignment between voxel data and coordinates.

\subsection*{Architectural Foundations: The Physical Coordinate System}
To address Challenge 4, software must rigorously maintain the mapping between discrete voxel indices $\mathbf{v}$ and physical patient space $\mathbf{p}$ via a $4 \times 4$ homogeneous affine transformation matrix $M$:
\begin{equation}
\mathbf{p} = M \mathbf{v}, \quad M = 
\begin{bmatrix} 
R S & T \\
\mathbf{0} & 1
\end{bmatrix}
\end{equation}
where $R$ is rotation, $S$ is scaling (spacing), and $T$ is translation (origin). MedImages.jl ensures that any spatial operation $\mathcal{T}$ applied to the voxel grid is coupled with an update $M_{new} = M_{old} \times M_{\mathcal{T}}$, preventing desynchronization.

\section{Methods}

\subsection*{Software Architecture: Zero-Cost Metadata Abstractions}
MedImages.jl implements a BIDS-inspired metadata paradigm where the \texttt{MedImage} struct enforces metadata integrity through Julia's type system. Unlike Python workflows where converting to arrays discards spatial context, MedImages.jl uses zero-cost abstractions—metadata (origin, spacing, direction, patient ID, acquisition timestamp) is intrinsic to the structure without runtime overhead.

ITK (C++) is used primarily for I/O robustness (NIfTI, DICOM parsing), while all transformations are implemented natively in Julia. This ensures end-to-end differentiability via Zygote.jl and Enzyme.jl. To support high-throughput multimodal studies, \texttt{BatchedMedImage} stores 4D voxel data $(x, y, z, \text{batch})$ with vectorized metadata, facilitating efficient GPU operations.

\subsection*{Experimental Setup: Evaluating the Four Challenges}

\subsubsection*{Evaluating Volume and Speed (Challenges 1 and 2)}
We developed a fused affine preprocessing pipeline that collapses orientation normalization, resampling, and spatial centering into a single matrix. We benchmarked a 100-case preprocessing pipeline on the autoPET dataset against MONAI's \texttt{PersistentDataset}.

\subsubsection*{Differentiability and Physics Discovery: The UDE Dosimetry Framework}
To validate Challenge 3, we implemented a 4-state UDE system to model $^{177}$Lu-PSMA dosimetry over a 300-hour therapeutic window. The system integrates biological activity kinetics with a learned radiation transport residual:
\begin{align}
    \frac{dA_{blood}}{dt} &= -(k_{10} + \lambda_{phys}) A_{blood}(t) - \sum \text{Uptake} + \sum \text{Washout} \\
    \dot{D}_{total}(\mathbf{r}, t) &= \text{Softplus} \left( \frac{A_{total}(\mathbf{r}, t) \cdot \lambda_{conv}}{m(\mathbf{r})} + \mathcal{N}_\theta(\mathbf{A}_{std}, \rho_{std}) \right)
\end{align}
The absorbed dose rate $\dot{D}$ is modeled as a sum of local energy deposition ($\lambda_{conv} = 8.478 \times 10^{-8}$) and a learned residual $\mathcal{N}_\theta$ (3D Residual Net) that discovers heterogeneous scattering physics from gold-standard Monte Carlo (MC) ground truth.

\subsubsection*{Fidelity Benchmarking and Baselines}
We contextualized the UDE performance ($N=48$ patches, 12 patients) against seven baselines:
\begin{enumerate}
    \item \textbf{Analytical Baseline (VSV):} Derived from the \texttt{Evaluate-Proj-Statistics-Dosimetry} framework. It utilizes monochromatic dose kernels and $^{177}$Lu sensitivity of \textbf{9.26 CPS/MBq}.
    \item \textbf{Monte Carlo Reference (MC-GT):} Gold-standard labels generated via OpenGATE/Geant4 using $10^9$ histories.
    \item \textbf{Stabilized CNN:} A ResNet-32 architecture mapping SPECT/CT to MC dose via GroupNorm and Huber loss.
    \item \textbf{PBPK-ML:} Extra Trees regressor using 12 local statistical features.
    \item \textbf{Spect0:} A residual learning U-Net designed to predict spatial corrections for the VSV baseline.
    \item \textbf{SemiDose (3D Mean Teacher):} A semi-supervised framework using consistency regularization between student and teacher networks to leverage unlabeled volumetric data.
    \item \textbf{DblurDoseNet:} A deep 3D U-Net mapping low-resolution SPECT and anatomical CT directly to high-fidelity dose maps via direct deconvolution learning.
\end{enumerate}

\section{Results}

\subsection*{Challenge 1 and 2: Performance Scaling and GPU Acceleration}
MedImages.jl achieved a \textbf{7.2$\times$ total turnaround speedup} over MONAI in high-volume biobank processing (Table \ref{tab:scaling_comparison}). Native GPU spatial kernels reached up to \textbf{135$\times$ speedups} over SimpleITK CPU baselines.

\begin{table}[H]
\centering
\caption{Performance Scaling Comparison (100 Subjects, Epochs 2+)}
\begin{tabular}{|l|c|c|c|}
\hline
\textbf{Phase (per Subject)} & \textbf{MedImages (HDF5)} & \textbf{MONAI (Cache)} & \textbf{Speedup} \\
\hline
Stage 1: Load from Cache & \textbf{~40 ms} & ~150 ms & 3.7$\times$ \\
Stage 2: Transform (Compute) & \textbf{~50 ms} & ~500 ms & 10.0$\times$ \\
\hline
\textbf{Total Turnaround} & \textbf{~90 ms} & \textbf{~650 ms} & \textbf{7.2$\times$} \\
\hline
\end{tabular}
\label{tab:scaling_comparison}
\end{table}

\subsection*{Challenge 3: Differentiability and SciML Dominance}
The Julia/SciML stack provided a \textbf{10.5--13.6$\times$ computational speedup} over equivalent PyTorch and JAX implementations (Fig \ref{fig:speed_comparison}). The UDE model achieved superior fidelity ($r=0.957$) by discovering scattering physics that clinical analytical models ignore.

\begin{figure}[H]
    \centering
    \includegraphics[width=0.7\textwidth]{elsarticle/figures_new/speed_comparison.png}
    \caption{\textbf{Cross-Language Performance Benchmark.} Julia's native JIT fusion provides an order-of-magnitude speed advantage for volumetric dosimetry ODEs.}
    \label{fig:speed_comparison}
\end{figure}

\begin{table}[H]
\centering
\caption{Dosimetry Performance Benchmark ($64^3$ Patches, $N=48$)}
\begin{tabular}{|l|c|c|c|}
\hline
\textbf{Model Architecture} & \textbf{Pearson $r$} & \textbf{MAE (norm)} & \textbf{Consistency} \\ \hline
\textbf{UDE (No-Approx)} & \textbf{0.957} & 0.033 & SciML-Integrated \\
Stabilized CNN (ResNet-32) & 0.939 & \textbf{0.027} & Image-to-Image \\
Analytical Baseline (VSV) & 0.912 & 0.032 & Physics-Only \\
PBPK-ML (Extra Trees) & 0.772 & 0.075 & Statistical ML \\
\hline
\end{tabular}
\label{tab:dosimetry_all}
\end{table}

\begin{figure}[H]
    \centering
    \includegraphics[width=\textwidth]{elsarticle/figures_new/dosimetry_comparison_all.png}
    \caption{\textbf{Comprehensive 2-row Dosimetry Comparison.} (A) Anatomy (CT). (B) Monte Carlo Ground Truth. (C) UDE Champion. (D) Stabilized CNN. (E) VSV Baseline. (F) PBPK-ML. (G) DblurDoseNet. (H) SemiDose. The UDE approach (C) provides the tightest structural and quantitative alignment with the MC physics.}
    \label{fig:dosimetry_all}
\end{figure}

\subsection*{Challenge 4: Metadata Integrity and Clinical Fidelity}
Quantitative consistency was maintained throughout compounded transformations, with $<1.5\%$ deviation in SUV and lesion volume (Fig \ref{fig:trans_consistency}).

\begin{figure}[H]
    \centering
    \includegraphics[width=0.8\textwidth]{elsarticle/figures_new/transformation_consistency.png}
    \caption{\textbf{Systematic Transformation Consistency Analysis.} MedImages.jl ensures absolute metadata integrity through zero-cost metadata coupling.}
    \label{fig:trans_consistency}
\end{figure}

\section{Discussion and Conclusions}
MedImages.jl represents a fundamental shift in medical image analysis. By bridging the divide between acquisition physics and inference logic, it empowers clinicians to utilize high-fidelity SciML models that are order-of-magnitude faster than hybrid frameworks. The success of the UDE framework ($r=0.957$) proves that native high-performance ecosystems are the future of quantitative theranostics.

\bibliographystyle{elsarticle-num}
\bibliography{bibl}

\end{document}
